\label{chap:conclusion}

This thesis set out to design, build and evaluate a hybrid graph-table interface for programme-level study planning, and to learn from students how such an interface is used. It did so in three steps that answer its three research questions: a formative study that produced a specification of what students need (Chapter~\ref{chap:formative-study}), a design and an implementation that realise it (Chapters~\ref{chap:design} and~\ref{chap:implementation}), and an evaluation that observed eleven students planning with the result, once without the curriculum graph and once with it (Chapter~\ref{chap:evaluation}), interpreted in Chapter~\ref{chap:discussion}. This chapter states the answers, restates what the thesis contributes, and names the work that follows.

\section{Answers to the Research Questions}
\label{sec:concl-rq-answers}

This section states what each research question was answered with, and how completely.

\paragraph{Research Question 1: students' needs and the resulting requirements for a hybrid graph-table interface supporting programme-level planning and the understanding of curriculum structure and course dependencies.}
The answer is the specification of Section~\ref{chap:from-themes-to-requirements}: 13 features carrying 55 numbered requirements, each traceable through a theme to the interview excerpts that motivated it, and ranked by an importance score that governed what was built first. It is answered on the evidence of six students at one faculty, which bounds it in the way Chapter~\ref{chap:methodology} states: the specification documents what this population needed, not what students in general need.

\paragraph{Research Question 2: how a hybrid interface combining graph-based visualisation with table-based planning can be designed and implemented from those requirements.}
The answer is the artefact together with the reasoning recorded alongside it. Chapter~\ref{chap:design} answers the design half by deriving each decision from a requirement, a prototype assumption, published guidance, or a stated design judgement, and by naming the alternative it was weighed against where one was live; Chapter~\ref{chap:implementation} answers the implementation half with a working system whose central structural decision, a single authoritative plan model from which both views and all feedback are derived, is what makes the two views projections of one plan rather than two plans to reconcile. The question asks \emph{how} such an interface can be built, and a single realisation cannot establish that this is the only way or the best way to build one; what it describes is one way it can be built, and at what cost in specific decisions.

\paragraph{Research Question 3: how students plan with the hybrid interface, and what the graph-based view contributes to their planning.}
Read through the lenses introduced in Chapter~\ref{chap:introduction}, the answer holds as follows:

\begin{itemize}
  \item \textbf{Process (The Core Planning Loop):} Students plan in a single, checklist-driven loop that was consistent across every participant. The student reads the Dashboard's missing-requirements list, finds and holds a candidate, places it, and then re-reads the list. The graph does not add a new station to this cycle; rather, its value sits in the query step (largely through its filters), where it substitutes for the catalogue. It does not support the commit step.

  \item \textbf{Process (Scaffolding and Staging):} The components sitting outside this core loop function as cognitive buffers and temporary scaffolding rather than primary drivers. A pre-filled backbone successfully bypasses the ``cold start'' problem, while a staging area (the Parking Stage) serves as a vital hinge, allowing users to decouple finding a course from committing it to a semester. Similarly, recommendations act as training wheels: while highly utilized initially, participants discarded them in favour of manual filtering once they could judge the curriculum themselves. Ultimately, students rely heavily on these peripheral tools to manage cognitive load, but when the system requires them to manually configure their own guardrails (such as workload thresholds), it creates silent failure states where unconfigured defaults go unnoticed. 

  \item \textbf{Plans Produced:} How students plan is fundamentally shaped by an over-reliance on the tool's authoritative signals. An independent audit of the twenty-one readable runs showed that while seventeen met the constraints explicitly modelled by the tool, only thirteen met the full task brief. Of the eight that failed, four missed the reduced-load requirement, the single task constraint the system's checklist could not track. Crucially, seven of those failed runs were rated as a complete success by the participants who produced them. Students offloaded their compliance tracking entirely to the interface, equating a clear dashboard with absolute success and letting unmodelled constraints fall out of their working memory.

  \item \textbf{Satisfaction:} Self-reported satisfaction was high and consistent, scoring higher on pragmatic quality than hedonic quality. Perspicuity (clarity) was the lowest-scoring scale, which aligns directly with the terminology and learnability struggles observed in the recordings.
\end{itemize}

\section{Future Work}
\label{sec:concl-future}

The implemented system provides a platform for future research. While this thesis focused on the planning process within a controlled lab setting, the tool could easily support longitudinal studies during actual enrolment periods, and comparisons against existing institutional tools with larger, non-computing cohorts. Building directly on the findings of this study, future work should pursue six specific directions:

\begin{itemize}
  \item \textbf{Further testing the value of the graph.} The study yields a falsifiable hypothesis: if the graph's value lies primarily in its filter panel, moving those filters into the tabular view should render the graph largely optional. A future study should test this by adding filters to the table and measuring whether the query step remains there.
  
  \item \textbf{Resolving critical UI bottlenecks.} The problem register (Table~\ref{tab:problems}) establishes a clear development agenda. For instance adding semester workload totals directly inside the graph would finally allow the graph to support the ``commit'' step, enabling researchers to test how users commit from a structural view.
  
  \item \textbf{Scaling up via telemetry.} The qualitative process account presented here should be scaled into a quantitative log study. Because the system already timestamps every state change (Chapter~\ref{chap:implementation}, Section~\ref{sec:impl-frontend}), it can seamlessly capture exact telemetry, dwell times, placements, and rejections, without manual coding. This would allow the descriptive loops identified here to be tested statistically on hundreds of students planning their actual degrees over a semester.
  
  \item \textbf{Visualising eligibility gates rather than sequences.} Four participants asked for ordering information at the point of placement, yet the tested curricula encode almost no hard sequential dependencies. Future research should investigate curricula with stricter ordering constraints, or reconsider what a curriculum graph ought to display in their absence. Rather than forcing a visual representation of a sparse prerequisite network, it may be significantly more useful to visualise broad eligibility gates (such as StEOP or core-before-elective rules).
  
  \item \textbf{Investigating expertise-moderated tool use.} The exploratory patterns observed here, particularly how domain expertise dictates whether a student uses automated recommendations or manual filters, warrant a study. This requires an explicit measure of prior knowledge and a broader sample drawn from outside computing programmes.
\end{itemize}